% =====================================================================
%  thmsmith Stage -1 proposal skeleton.
%  Written automatically by the thmsmith TS pipeline before Stage -1 runs.
%  Locked parameters: qid=pid\_ordinal\_principal\_facet\_phase, specialization=v1
% =====================================================================

\documentclass[11pt]{article}

\usepackage{amsmath,amssymb,amsthm,mathtools}
\usepackage{enumitem}
\usepackage[colorlinks=true,linkcolor=blue,citecolor=blue,urlcolor=blue]{hyperref}
\usepackage{natbib}

\newcommand{\E}{\mathbb{E}}
\renewcommand{\Pr}{\mathbb{P}}
\newcommand{\one}{\mathbf{1}}
\newcommand{\transpose}{\mathsf{T}}
\newcommand{\calH}{\mathcal{H}}
\newcommand{\calF}{\mathcal{F}}
\newcommand{\calR}{\mathcal{R}}
\newcommand{\R}{\mathbb{R}}
\newcommand{\plim}{\operatorname*{plim}}

\theoremstyle{definition}
\newtheorem{definition}{Definition}
\newtheorem{assumption}{Assumption}
\newtheorem{example}{Example}

\theoremstyle{plain}
\newtheorem{theorem}{Theorem}
\newtheorem{conjecture}{Conjecture}
\newtheorem{proposition}{Proposition}
\newtheorem{lemma}{Lemma}
\newtheorem{corollary}{Corollary}

\theoremstyle{remark}
\newtheorem{remark}{Remark}

\title{thmsmith Stage -1 proposal: \texttt{pid\_ordinal\_principal\_facet\_phase} / \texttt{v1}%
  \\ \large\textit{K=3 middle-stratum compatibility facet and phase threshold for ordinal principal-strata bounds}}
\author{thmsmith Stage -1 proposer}
\date{\today}

\begin{document}
\maketitle

% --- BEGIN CONTENT ---  (everything above is fixed by the pipeline)

\paragraph{Locked parameters.}
\begin{verbatim}
qid=pid_ordinal_principal_facet_phase
specialization=v1
anchor_cluster=Partial identification
observable_distribution_P=law of (Z,S,Y,X), where Z in {0,1} is randomized,
  S in {0,1,2} for the flagship K=3 cell and may lie in {0,...,K-1}
  in reductions, Y in finite support Ygrid subset [0,1], and X in finite
  support Xgrid
parameter_theta=theta_g = mu_1g - mu_0g for a fixed principal stratum
  g=(s0,s1), with mu_zg=E[Y(z)|G=g] and G=(S(0),S(1))
identifying_restrictions=randomization, consistency, positivity, supplied
  covariate-conditional ordinal principal scores pi_g(x), simplex and
  principal-score calibration of posterior stratum weights q_{z,g}(x,y), and
  Gamma-adjacent odds-ratio divergence inside each observed (Z=z,S=k,X=x,Y=y)
  cell
bounding_functional=the common-law feasible polytope Q(P,pi,Gamma), its
  K=3 middle-facet certificate C_mid(P,pi,Gamma), the ordinal sharp interval
  [L_g(P,pi,Gamma),U_g(P,pi,Gamma)], and the symmetric witness slack
  Delta_mid_sym(Gamma) that certifies the non-PI phase Gamma>1 by comparing
  two active-basis endpoint formulas derived from the same K=3 cell primitives
\end{verbatim}

\begin{abstract}
Ordinal principal-strata sensitivity analysis has a local obstruction that is invisible in adjacent two-stratum reductions.
The parent result \texttt{pid\_ordinal\_principal\_strata/v1} proves a sharp adjacent-chain posterior LP and records nonrectangularity, but it leaves the mechanism bundled inside generic feasibility.
This upgrade isolates the K=3 shared-middle mechanism as a named compatibility facet with an explicit normal, an endpoint-attaining rational witness law, and a primitive phase slack for when the facet controls the bound beyond all pairwise adjacent relaxations.
The main statements are conjectural at flagship level: the facet and phase claims must be checked as irredundant polyhedral statements, not as a restatement of LP duality.
If resolved, the result would give ordinal principal-strata sensitivity analyses a hand-checkable certificate for when dichotomizing an ordinal intercurrent event loses sharp information.
\end{abstract}

\section{Introduction}
Principal strata defined by an ordinal post-treatment variable share latent middle strata across adjacent reductions.  The main claim is that, already for K=3, the shared middle stratum yields a specific compatibility facet whose endpoint separates the common-law feasible polytope from the intersection of all pairwise adjacent-stratum relaxations.

Relative to the parent qid \texttt{pid\_ordinal\_principal\_strata/v1}, the upgrade delta is mechanical: the parent has sharpness and a strict-gap witness, while this run names the facet certificate \(C_{\mathrm{mid}}\), computes its normal, and conjectures the symmetric witness slack \(\Delta_{\mathrm{mid}}^{\mathrm{sym}}\) that marks the non-PI phase \(\Gamma>1\).  Relative to \citet[p.~21]{FrangakisRubin2002} and \citet[Section~2]{DingLu2017}, the result keeps the common principal-strata law and principal-score mixture structure but studies a fixed-pi partial-identification geometry outside principal ignorability.  The kernel lives in partial identification.

The concrete consumer is ordinal intercurrent-event principal-stratum analysis in clinical trials, including the sensitivity workflows represented by \citet{WangZhangMealliBornkamp2023} and the margin-free odds-ratio principal-stratification framework in \citet{TongKahanHarhayLi2025}.  A reader who accepts the conclusion should report the middle-facet certificate, or explain why the ordinal analysis was replaced by pairwise binary sensitivity bounds.

\begin{itemize}[leftmargin=*]
\item Theorem 1: the parent adjacent-chain sharp set reduces to the common-law feasible polytope used here. Gap: \texttt{pid\_ordinal\_principal\_strata/v1} has the LP but not the mechanism certificate. Consumer: ordinal principal-score sensitivity analysis.
\item Conjecture 1: the K=3 middle-facet certificate \(C_{\mathrm{mid}}\) is a valid and irredundant facet of the common-law cell polytope. Gap: \citet[Section~2, Theorem~1]{LuZhangDing2020} gives endpoint constructions for ordinal treatment-effect bounds from potential-outcome marginals, not a principal-strata shared-middle posterior facet. Consumer: finite-support ordinal principal-strata bounds.
\item Conjecture 2: in the symmetric K=3 witness cell, the primitive slack \(\Delta_{\mathrm{mid}}^{\mathrm{sym}}(\Gamma)=(\Gamma+2)/\{3(\Gamma+1)(2\Gamma+1)\}\) is positive exactly for the non-PI phase \(\Gamma>1\), and the edgewise endpoint cannot be embedded in a common law. Gap: \citet[Section~3.2]{DingLu2017} identifies under principal ignorability, and \citet[Sections~2 and 5]{TongKahanHarhayLi2025} defines odds-ratio principal-stratification sensitivity without a K=3 ordinal compatibility phase. Consumer: clinical-trial intercurrent-event sensitivity workflows.
\end{itemize}

\section{Related work}
\citet[p.~21]{FrangakisRubin2002} [reused from parent] define principal strata by joint potential values of a post-treatment variable; this proposal keeps that common latent law as the object whose adjacent projections must be compatible.
\citet[Section~2]{DingLu2017} [reused from parent] identify principal causal effects using principal scores under principal ignorability; the present question begins once the principal-score posterior is allowed to tilt with the outcome.
\citet{JiangYangDing2022} [reused from parent] develop multiply robust estimation under principal ignorability, which is a point-identified baseline rather than a fixed-pi partial-identification geometry.
\citet{FellerMealliMiratrix2017} emphasize that principal ignorability is strong and hard to justify in practice, motivating a bound whose tightening comes from compatibility rather than from stronger ignorability.

\citet{LuDingDasgupta2018} study ordinal-outcome treatment-benefit estimands whose sharp bounds depend on unobserved joint potential-outcome distributions; their relevance is the ordinal sharp-bound frontier, while the present latent object is the posterior principal-strata composition inside an observed cell.
\citet[Section~2, Theorem~1]{LuZhangDing2020} give explicit sharp bounds and endpoint-attaining constructions for ordinal relative treatment effects; this proposal imports the endpoint-certificate standard but applies it to a different K=3 principal-strata compatibility facet.
\citet{ZhangRubin2003} [reused from Stage -1.1] show that principal-effect bounds often require explicit latent-stratum constructions, which is the proof style required by the K=3 witness law here.
\citet{VanderWeele2011} [reused from parent] warns that dichotomizing ordinal or continuous intermediates can mislead; the pairwise-relaxation comparator formalizes that warning in one finite-support principal-strata cell.

\citet{WangZhangMealliBornkamp2023} [reused from parent] give a multiple-imputation sensitivity workflow for principal ignorability in a binary intercurrent-event setting; the proposed facet says what must be added when the intercurrent event has three ordered levels.
\citet[Sections~2 and 5]{TongKahanHarhayLi2025} is a new frontier anchor: it uses a conditional odds-ratio sensitivity parameter for semiparametric principal stratification beyond monotonicity and motivates a nonparametric finite-cell certificate when the principal stratum is ordinal.
\citet{WuAntonelli2024} [reused from parent] studies partial identification under principal-ignorability violations, but does not isolate a K=3 adjacent posterior-chain facet.
\citet{JiangDing2021} marks the boundary between auxiliary-variable identification and fixed-pi compatibility geometry.

The in-repository parent \texttt{pid\_ordinal\_principal\_strata/v1} proves sharp finite-support adjacent-chain bounds and a strict edgewise gap; its accepted tier was field because the nonrectangularity mechanism remained bundled inside the LP.
The live \texttt{doc/API.md} PartialID catalogue records no promoted ordinal-principal entry beyond seed-level notes and a Manski study wrapper, so the duplicate risk is the parent itself.
This upgrade is not a new-domain LP dual: its claimed delta is the named middle-facet normal, endpoint-attaining K=3 witness, and primitive phase slack.

\section{Formal setup}
Let \(P\) be the law of \((Z,S,Y,X)\), with randomized \(Z\in\{0,1\}\), ordinal \(S\in\{0,1,2\}\) in the flagship cell, finite \(\mathcal Y\subset[0,1]\), and finite \(\mathcal X\).  Potential intermediates are \(S(0),S(1)\), the principal stratum is \(G=(S(0),S(1))\), and \(\pi_g(x)=\Pr(G=g\mid X=x)\) is supplied.  For an observed arm-level cell \((z,k,x)\), define \(A_z(k)=\{h:h_z=k\}\), baseline cell weights \(w_h^z(x)=\pi_h(x)/\sum_{\ell\in A_z(k)}\pi_\ell(x)\), and posterior weights
\[
  q_{z,h}(x,y)=\Pr(G=h\mid Z=z,S=k,X=x,Y=y),\qquad h\in A_z(k).
\]
The common-law feasible set \(\mathcal Q(P,\pi,\Gamma)\) consists of posterior arrays satisfying simplex constraints, calibration to \(w^z(x)\), and adjacent odds-ratio inequalities
\[
  \Gamma^{-1}q_{z,\ell}(x,y)w_h^z(x)
  \le q_{z,h}(x,y)w_\ell^z(x)
  \le \Gamma q_{z,\ell}(x,y)w_h^z(x)
\]
for every adjacent pair \(h\sim_z\ell\).  The target is \(\theta_g=\mu_{1g}-\mu_{0g}\), where
\[
  \mu_{zg}(q)=\frac{\E_P[\one\{Z=z,S=g_z\}q_{z,g}(X,Y)Y/e_z]}{p_g},
  \qquad p_g=\E\{\pi_g(X)\}.
\]
The ordinal sharp interval is \([L_g,U_g]=[\min_{q\in\mathcal Q}\{\mu_{1g}(q)-\mu_{0g}(q)\},\max_{q\in\mathcal Q}\{\mu_{1g}(q)-\mu_{0g}(q)\}]\).

For the K=3 local certificate, fix one positive cell \((z,k,x_0)\) with three adjacent strata \(a\sim_z b\sim_z c\), outcome support \(\{0,1\}\), and let \(\lambda=\Pr(Y=1\mid Z=z,S=k,X=x_0)\).  Write \(q_h^1=q_{z,h}(x_0,1)\), \(q_h^0=q_{z,h}(x_0,0)\), and \(w_h=w_h^z(x_0)\).  Under the symmetric witness submodel \(\lambda=1/2\) and \(w_a=w_b=w_c=1/3\), the middle-facet endpoint is
\[
  T_{\mathrm{mid}}(\Gamma)=\frac{4\Gamma-1}{3(2\Gamma+1)}.
\]
The middle-facet certificate is the inequality
\[
  C_{\mathrm{mid}}(q;\Gamma):\quad q_b^1\le T_{\mathrm{mid}}(\Gamma),
\]
with active normal \(n_{\mathrm{mid}}=e_b^1\) in the projected endpoint coordinate and derivation certificate
\[
  (q_a^0-\Gamma q_b^0)+(q_c^0-\Gamma q_b^0)\le0,
  \qquad q_a^0+q_b^0+q_c^0=1,
  \qquad q_b^0+q_b^1=2/3.
\]
For \(\Gamma>1\), the pairwise adjacent relaxation endpoint for one incident edge in the same symmetric cell is
\[
  T_{\mathrm{pair}}(\Gamma)=\frac{2\Gamma+1}{3(\Gamma+1)}.
\]
The symmetric primitive slack is
\[
  \Delta_{\mathrm{mid}}^{\mathrm{sym}}(\Gamma)
  =
  (\Gamma+2)\{3(\Gamma+1)(2\Gamma+1)\}^{-1}.
\]
It is obtained by substituting the same cell primitives into the coherent active-basis certificate and the one-edge active-basis certificate, not by defining the phase as an endpoint difference.  For \(\Gamma>1\), the algebraic identity
\[
  T_{\mathrm{pair}}(\Gamma)-T_{\mathrm{mid}}(\Gamma)
  =
  \Delta_{\mathrm{mid}}^{\mathrm{sym}}(\Gamma)
\]
is the witness-phase claim.  At \(\Gamma=1\), the maintained common-law model collapses to the principal-score posterior \(q=w\); the edgewise endpoint formula above is not used as a PI-continuous boundary formula because its active basis changes when the odds-ratio inequalities become equalities.

\begin{verbatim}
qid=pid_ordinal_principal_facet_phase
specialization=v1
anchor_cluster=Partial identification
observable_distribution_P=law of (Z,S,Y,X), where Z in {0,1} is randomized,
  S in {0,1,2} for the flagship K=3 cell and may lie in {0,...,K-1}
  in reductions, Y in finite support Ygrid subset [0,1], and X in finite
  support Xgrid
parameter_theta=theta_g = mu_1g - mu_0g for a fixed principal stratum
  g=(s0,s1), with mu_zg=E[Y(z)|G=g] and G=(S(0),S(1))
identifying_restrictions=randomization, consistency, positivity, supplied
  covariate-conditional ordinal principal scores pi_g(x), simplex and
  principal-score calibration of posterior stratum weights q_{z,g}(x,y), and
  Gamma-adjacent odds-ratio divergence inside each observed (Z=z,S=k,X=x,Y=y)
  cell
bounding_functional=the common-law feasible polytope Q(P,pi,Gamma), its
  K=3 middle-facet certificate C_mid(P,pi,Gamma), the ordinal sharp interval
  [L_g(P,pi,Gamma),U_g(P,pi,Gamma)], and the symmetric witness slack
  Delta_mid_sym(Gamma) that certifies the non-PI phase Gamma>1 by comparing
  two active-basis endpoint formulas derived from the same K=3 cell primitives
\end{verbatim}

\section{Assumptions}
\begin{assumption}[Randomization and consistency]\label{ass:randomization}
\((Y(0),Y(1),S(0),S(1),X)\perp Z\), \(S=S(Z)\), and \(Y=Y(Z)\).
\end{assumption}

\begin{assumption}[Finite support and positivity]\label{ass:positivity}
\(\mathcal X\), \(\mathcal Y\), and the principal-strata set are finite, \(\mathcal Y\subset[0,1]\), \(0<e_z<1\), and \(p_g=\E\{\pi_g(X)\}>0\) for the target stratum.  All displayed conditional probabilities and odds ratios are evaluated only on positive-probability observed cells.
\end{assumption}

\begin{assumption}[Supplied principal scores]\label{ass:principal-scores}
For each \(x\), \(\pi(x)\) is a probability vector on principal strata and is compatible with the observed intermediate margins:
\[
  \Pr(S=k\mid Z=z,X=x)=\sum_{h\in A_z(k)}\pi_h(x).
\]
\end{assumption}

\begin{assumption}[Adjacent posterior divergence]\label{ass:divergence}
A posterior array \(q\) is feasible when it satisfies simplex and calibration in each observed cell and the adjacent odds-ratio inequalities in Section~6 for a fixed \(\Gamma\in[1,\infty]\).
\end{assumption}

\begin{assumption}[K=3 symmetric witness cell]\label{ass:witness-cell}
For the flagship certificate, there is one positive cell \((z,k,x_0)\) with three adjacent strata \(a\sim_z b\sim_z c\), \(\mathcal Y=\{0,1\}\) on that cell, \(\Pr(Y=1\mid Z=z,S=k,X=x_0)=1/2\), \(w_a=w_b=w_c=1/3\), and \(1<\Gamma<\infty\).
\end{assumption}

\begin{assumption}[Nondegenerate phase neighborhood]\label{ass:phase-neighborhood}
The local cell has positive slack in all nonactive simplex, calibration, positivity, and margin-compatibility constraints, so small relative perturbations of \((P,\pi)\) preserve the active basis that defines \(T_{\mathrm{mid}}\) and \(T_{\mathrm{pair}}\).
\end{assumption}

\section{Main results}
\begin{theorem}[Theorem 1: parent sharp set as common-law polytope]\label{thm:parent-reduction}
Under Assumptions~\ref{ass:randomization}--\ref{ass:divergence}, the sharp set from \texttt{pid\_ordinal\_principal\_strata/v1} is the image of \(\mathcal Q(P,\pi,\Gamma)\) under the linear map \(q\mapsto\mu_{1g}(q)-\mu_{0g}(q)\).  Thus the present proposal reduces to the parent when the only reported object is \([L_g(P,\pi,\Gamma),U_g(P,\pi,\Gamma)]\).
\end{theorem}
The proof is the parent Bayes-disintegration argument: every latent law induces a calibrated posterior array, and every feasible posterior array can be completed to a latent law with observable margin \(P\).

\begin{conjecture}[Conjecture 1: K=3 middle-facet certificate (NOT YET PROVED)]\label{conj:middle-facet}
Under Assumptions~\ref{ass:randomization}--\ref{ass:witness-cell}, the certificate \(C_{\mathrm{mid}}(q;\Gamma):q_b^1\le T_{\mathrm{mid}}(\Gamma)\) is a valid and irredundant facet of the projected common-law cell polytope for the endpoint coordinate \(q_b^1\).  The facet is supported by the normal \(n_{\mathrm{mid}}=e_b^1\), and its active face is attained by
\[
  q_b^0=\frac{1}{2\Gamma+1},\quad
  q_a^0=q_c^0=\frac{\Gamma}{2\Gamma+1},\quad
  q_b^1=\frac{4\Gamma-1}{3(2\Gamma+1)},\quad
  q_a^1=q_c^1=\frac{\Gamma+2}{3(2\Gamma+1)}.
\]
Moreover, for every finite \(\Gamma>1\), no pair of independent adjacent-edge optimizers attaining \(T_{\mathrm{pair}}(\Gamma)\) on both incident edges has a common middle coordinate in the symmetric witness cell, because the middle coordinate required by each edge exceeds the coherent facet endpoint by \(\Delta_{\mathrm{mid}}^{\mathrm{sym}}(\Gamma)>0\).
\end{conjecture}

\begin{conjecture}[Conjecture 2: middle-facet phase threshold (NOT YET PROVED)]\label{conj:phase}
Under Assumptions~\ref{ass:randomization}--\ref{ass:phase-neighborhood}, the symmetric K=3 witness has a non-PI strict-containment phase exactly for \(\Gamma>1\).  In that cell,
\[
  \Delta_{\mathrm{mid}}^{\mathrm{sym}}(\Gamma)
  =\frac{\Gamma+2}{3(\Gamma+1)(2\Gamma+1)}>0,
\]
and the active-basis identities imply
\[
  T_{\mathrm{pair}}(\Gamma)-T_{\mathrm{mid}}(\Gamma)
  =
  \Delta_{\mathrm{mid}}^{\mathrm{sym}}(\Gamma).
\]
Thus the incident pairwise endpoint is separated from the coherent common-law endpoint by a primitive slack computed from the K=3 cell, not by a phase variable defined as the endpoint difference.  The conjectural extension is that the same active basis remains valid on a relative-open neighborhood of the symmetric witness inside the margin-compatible parameter space.
\end{conjecture}

\section{Sanity-check corollaries}
\begin{corollary}[Reduction to the parent]\label{cor:parent}
Under Assumptions~\ref{ass:randomization}--\ref{ass:divergence}, if \(C_{\mathrm{mid}}\) and \(\Delta_{\mathrm{mid}}^{\mathrm{sym}}\) are not reported, Theorem~\ref{thm:parent-reduction} gives exactly the parent sharp interval \([L_g,U_g]\).  The upgrade is therefore a mechanism refinement of the parent rather than a different estimand.
\end{corollary}

\begin{corollary}[Binary and PI corners]\label{cor:corners}
If \(K=2\), no cell contains a shared middle stratum and \(C_{\mathrm{mid}}\) is absent.  If \(\Gamma=1\) on a connected positive-support cell, adjacent ratios force \(q=w\) and the principal-score formula is recovered; the symmetric edgewise formula for \(T_{\mathrm{pair}}\) is a \(\Gamma>1\) active-basis certificate and is not used for this PI corner.
\end{corollary}

\begin{example}[Exhibit 9.1: K=3 endpoint certificate]\label{ex:witness}
Set \(K=3\), \(X=\{x_0\}\), \(Y=\{0,1\}\), target \(b=(1,1)\), and in the treated cell \(Z=1,S=1\) let \(a=(0,1)\), \(b=(1,1)\), \(c=(2,1)\), \(w_a=w_b=w_c=1/3\), and \(\Pr(Y=1\mid Z=1,S=1,X=x_0)=1/2\).  For \(\Gamma=2\), the coherent endpoint is
\[
  (q_a^0,q_b^0,q_c^0,q_a^1,q_b^1,q_c^1)
  =\left(\frac25,\frac15,\frac25,\frac4{15},\frac7{15},\frac4{15}\right),
\]
so \(T_{\mathrm{mid}}(2)=7/15\).  The incident pairwise relaxation can attain
\[
  (r_a^0,r_b^0,r_a^1,r_b^1)=\left(\frac29,\frac19,\frac49,\frac59\right),
\]
so \(T_{\mathrm{pair}}(2)=5/9\) and \(\Delta_{\mathrm{mid}}^{\mathrm{sym}}(2)=4/45\).  The internal consistency checks are: all six coherent coordinates are in \([0,1]\), both outcome-state simplexes sum to one, each stratum calibration sums to \(2/3\), all denominators are positive, and the binary endpoint cannot be embedded because it would require \(q_b^1=5/9>7/15\).
\end{example}

\section{Proof-sketch route}
For Theorem~\ref{thm:parent-reduction}, the route is parent sharpness \(\leftarrow\) Bayes disintegration \(\leftarrow\) finite-support calibration.  For Conjecture~\ref{conj:middle-facet}, the route is: headline \(\leftarrow\) unfold the K=3 cell definitions \(\leftarrow\) substitute Assumptions~\ref{ass:witness-cell} and \ref{ass:divergence} \(\leftarrow\) derive the aggregate middle-deficit inequality from the two active adjacent inequalities at \(Y=0\) \(\leftarrow\) combine it with simplex and calibration to obtain \(q_b^1\le T_{\mathrm{mid}}\) \(\leftarrow\) verify affine independence of the active constraints on the projected face \(\leftarrow\) exhibit the rational endpoint in Exhibit~\ref{ex:witness}.  For Conjecture~\ref{conj:phase}, the route is: headline \(\leftarrow\) compute the two active-basis certificates from the same symmetric cell primitives \(\leftarrow\) simplify their algebraic slack to \(\Delta_{\mathrm{mid}}^{\mathrm{sym}}(\Gamma)\) \(\leftarrow\) verify \(\Delta_{\mathrm{mid}}^{\mathrm{sym}}(\Gamma)>0\) exactly for \(\Gamma>1\) \(\leftarrow\) show that the separated edge endpoint requires incompatible middle copies \(\leftarrow\) use active-basis stability of finite linear programs to promote the witness to a relative-open phase neighborhood.  Each conjectural flagship headline has at least four substantive steps beyond definition unfolding: aggregate-inequality derivation, endpoint construction, affine-independence or incompatibility verification, and active-basis stability.

\section{Honest scope flags}
Theorem~\ref{thm:parent-reduction} is inherited from the parent and is not the flagship novelty.  Conjectures~\ref{conj:middle-facet} and \ref{conj:phase} are open: the facet irredundancy, the exact active-basis phase claim, and the relative-open promotion still require proof.  The phase statement is deliberately limited to the symmetric K=3 witness cell; extending it to arbitrary unequal \((\lambda,w_a,w_b,w_c)\) may require a piecewise formula rather than one global expression.  The supplied-pi assumption is deliberate, because introducing principal-score uncertainty would confound the fixed-pi compatibility mechanism with a separate projection problem.

\section{Iteration trail}
\begin{itemize}[leftmargin=*]
\item v1 (cold-start) -- initial upgrade draft isolating the K=3 middle-facet certificate, endpoint witness, and phase scalar relative to the parent adjacent-chain LP; prior verdict: none.
\item v2 (revise) -- added location-level flagship anchors, replaced the tautological \(\Phi_{\mathrm{mid}}\) iff with the primitive symmetric slack \(\Delta_{\mathrm{mid}}^{\mathrm{sym}}\), limited the phase claim to the K=3 witness plus relative-open extension, and fixed the \(\Gamma=1\) corner; prior verdict: REVISE.
\end{itemize}

\section*{Bibliography}
\begin{thebibliography}{99}
\bibitem[Frangakis and Rubin(2002)]{FrangakisRubin2002}
Frangakis, C. E. and Rubin, D. B. (2002).
Principal stratification in causal inference.
\emph{Biometrics} 58, 21--29.

\bibitem[Ding and Lu(2017)]{DingLu2017}
Ding, P. and Lu, J. (2017).
Principal stratification analysis using principal scores.
\emph{Journal of the Royal Statistical Society: Series B} 79, 757--777.

\bibitem[Jiang, Yang, and Ding(2022)]{JiangYangDing2022}
Jiang, Z., Yang, S., and Ding, P. (2022).
Multiply robust estimation of causal effects under principal ignorability.
\emph{Journal of the Royal Statistical Society: Series B} 84, 1423--1445.

\bibitem[Feller, Mealli, and Miratrix(2017)]{FellerMealliMiratrix2017}
Feller, A., Mealli, F., and Miratrix, L. (2017).
Principal score methods: assumptions, extensions, and practical considerations.
\emph{Journal of Educational and Behavioral Statistics} 42, 726--758.

\bibitem[Lu, Ding, and Dasgupta(2018)]{LuDingDasgupta2018}
Lu, J., Ding, P., and Dasgupta, T. (2018).
Treatment effects on ordinal outcomes: causal estimands and sharp bounds.
\emph{Journal of Educational and Behavioral Statistics} 43, 540--567.

\bibitem[Lu, Zhang, and Ding(2020)]{LuZhangDing2020}
Lu, J., Zhang, Y., and Ding, P. (2020).
Sharp bounds on the relative treatment effect for ordinal outcomes.
\emph{Biometrics} 76, 664--669.

\bibitem[Zhang and Rubin(2003)]{ZhangRubin2003}
Zhang, J. L. and Rubin, D. B. (2003).
Estimation of causal effects via principal stratification when some outcomes are truncated by death.
\emph{Journal of Educational and Behavioral Statistics} 28, 353--368.

\bibitem[VanderWeele(2011)]{VanderWeele2011}
VanderWeele, T. J. (2011).
Principal stratification--uses and limitations.
\emph{International Journal of Biostatistics} 7, Article 28.

\bibitem[Wang et~al.(2023)]{WangZhangMealliBornkamp2023}
Wang, C., Zhang, Y., Mealli, F., and Bornkamp, B. (2023).
Sensitivity analyses for the principal ignorability assumption using multiple imputation.
\emph{Pharmaceutical Statistics} 22, 64--78.

\bibitem[Tong et~al.(2025)]{TongKahanHarhayLi2025}
Tong, J., Kahan, B., Harhay, M. O., and Li, F. (2025).
Semiparametric principal stratification analysis beyond monotonicity.
\emph{Statistica Sinica}, forthcoming.

\bibitem[Wu and Antonelli(2024)]{WuAntonelli2024}
Wu, M. and Antonelli, J. (2024).
Partial identification of principal causal effects under violations of principal ignorability.
\emph{arXiv:2412.06628}.

\bibitem[Jiang and Ding(2021)]{JiangDing2021}
Jiang, Z. and Ding, P. (2021).
Identification of causal effects within principal strata using auxiliary variables.
\emph{Statistical Science} 36, 493--508.
\end{thebibliography}

\end{document}
